\paragraph{An internal observer reads the protected response.}
The \texttt{protected\_particle\_readout} construction attaches the preceding
cube spectrum to complete native pointer records.  Its identified object is
the eight-dimensional cyclic response generated by the centered local density
probe, inside the 119-state $\eta=3$ sector and the larger 1820-dimensional
protected space.  It is a neutral, fixed-$N=6$ response; these three dimensions
are not interchangeable particle counts.  The construction uses Ramsey
interference to read a correlation function, a standard measurement principle,
and supplies its actual source implementation for this field.\footnote{For
many-body Ramsey correlation measurements, see Knap et al.,
\emph{Phys. Rev. Lett.} 111, 147205 (2013),
\url{https://arxiv.org/abs/1307.0006}.}

The apparatus retains sixteen field bits and one pointer, hence all
$2^{17}=131\,072$ field/pointer occupation rows.  The unchanged source prepares
$|R\rangle$ and a pointer $|+\rangle$, then applies controlled
$F_0=\exp(i\pi n_0/2)$.  There are two field segments of duration $1/2$:
the first is followed by complete C repair on all five admitted seams, and
the second by complete E repair on those seams.  Controlled $F_0^\dagger$,
a pointer phase $\phi\in\{0,\pi/2\}$, and the inverse preparation rotation
precede the complete six-outcome pointer export.  Every source first-failure
record and all 130 repair records remain present.  An outside repair outcome
continues; a native repair outcome or source failure terminates in its own
retained record.  No probability is conditioned on survival.

With exact field evolution, each repair acts identically on both ideal arms.
Immediately before the final pointer rotation their field vectors are
\[
 b_0=e^{-itH}|R\rangle,\qquad
 b_1=e^{i\phi}F_0^\dagger e^{-itH}F_0|R\rangle.
\]
Consequently the unconditional ideal pointer-zero probability is
\begin{align*}
 p_0(t,\phi)&=\tfrac12\{1+\operatorname{Re}[e^{i\phi}C(t)]\},\\
 C(t)&=\frac1{16}+\frac{15}{16}\sum_{j=1}^{8}
 w_j e^{-it\omega_j},\qquad \omega_j=z_j/16-1.
\end{align*}
Thus the two settings read $\operatorname{Re}C=2p_0(t,0)-1$ and
$\operatorname{Im}C=1-2p_0(t,\pi/2)$.  The phase $3t$ follows from the
reference eigenenergy; it is not an adjustable frequency fit.  A common
scalar added to $H$, or a common number chemical potential on this fixed
number domain, cancels between the arms.  Absolute energy zero is therefore
not determined by this experiment.

\paragraph{Complete source and field error, without a sector projection.}
The executor composes success and every first-export failure map of each
local word on at most three work bits, then applies these maps to the full
field/pointer array.  This is finite block contraction of the same native
word, not replacement by an ideal gate.  Independent tests compare the
contracted maps with direct source execution, including order-one failure
controls.  The exact $N=6$ restriction is used only by the independent
original-$H$ comparator, not by the native execution.

Each half-time field segment uses 32 symmetric product steps.  The actual
ordered support list has 24 hopping terms of norm $1/16$ and eight contact
terms of norm one.  Counting only possibly nonzero nested commutators gives
the all-Fock coefficient $K=629/256$.  The two segments therefore contribute
\[
 \varepsilon_{\rm split}
 =2\,\frac{629}{256}\frac{(1/2)^3}{32^2}
 =\frac{629}{1048576}.
\]
Independent rational trigonometric and literal-port bounds add less than
$8.56\times10^{-11}$ for the preparation, controlled pointer operations,
and both complete source field words.  The complete output-vector error
is consequently $\varepsilon<6.00\times10^{-4}$.  Every halted record is
part of this norm.  Complete repair and export preserve the comparison by
isometric extension or contraction; the all-input field bounds also hold
with untouched references.  The initial preparation bound consumes the
specified initialized blank.

For normalized complete outputs, half trace distance is at most their
vector distance.  Each unconditional pointer probability therefore differs
from its ideal value by at most $\varepsilon$, and each reconstructed
quadrature by at most $2\varepsilon$.  Rational root, residue and
trigonometric enclosures at $t=1$ imply, for the actual specified source,
\[
 0.993649<p_0(1,0)<0.994850,\qquad
 0.501379<p_0(1,\pi/2)<0.502580.
\]
Both intervals exclude the corresponding no-response values, one and one
half.  Numerical source residuals have their separate declared admission;
they are not substituted for these proved source bounds.
The executed construction uses 681\,283 native pulses per setting.
It gives pointer-zero probabilities approximately $0.9942497011$ and
$0.5019798894$, versus the independent original-$H$ values
$0.9942494663$ and $0.5019797169$.  The complete output-vector discrepancy
is approximately $2.27\times10^{-6}$, below the rigorous bound.

The imaginary fringe also distinguishes the contact interaction from the
same free cube.  Set $U=0$ while retaining the hopping, reference, local
probe and ideal repair schedule.  Bipartite occupation parity
$\Gamma_{\rm sub}=(-1)^{\sum_{x\,\mathrm{odd}}n_x}$ satisfies
\[
 \Gamma_{\rm sub}H_0\Gamma_{\rm sub}=-H_0,\qquad
 H_0|R\rangle=0,\qquad
 \Gamma_{\rm sub}F_0|R\rangle=F_0|R\rangle.
\]
The free characteristic function is consequently real at every time, so
$p_{0,\rm free}(t,\pi/2)=1/2$ exactly.  Both free arms remain in the
protected sector.  The certified interacting interval therefore excludes
this specified contact-free response.  This is discrimination between two
internal dynamics, not a fit of a physical particle mass.

For independent reset repetitions, Hoeffding's inequality and a union bound
give simultaneous sampling error at most $1/2048$ in these two probabilities
using $2^{24}$ trials per setting, except with probability less than
$4(3/8)^8$.\footnote{The bounded-independent-sum inequality is imported from
Hoeffding, \emph{J. Amer. Statist. Assoc.} 58 (1963), 13--30,
\url{https://doi.org/10.1080/01621459.1963.10500830}.}
The sampling allowance separates both fringes from no response.
Independent resets are the sampling-service premise; a seeded numerical
sample is an illustration of that model, not laboratory data.
The recorded PCG64 illustration, with seed 20261101, gives respectively
16\,680\,454 and 8\,423\,380 pointer-zero outcomes out of $2^{24}$ trials;
the complementary count includes every other complete outcome.

\paragraph{What eight-line resolution would additionally require.}
The executed two settings attach the known protected spectral measure to
an internally accessible record.  They do not reconstruct eight unknown
frequencies or individually resolve the eight lines.  There is, however,
a certified prospective experiment on the source-specified support.
At times $t_n=39n$, $n=0,\ldots,7$, use both pointer quadratures and set
\[
 y_n=\frac{16C(39n)-1}{15},\qquad
 V_{nj}=e^{-i39n\omega_j},\qquad y=Vw.
\]
Exact interval sine bounds in the Vandermonde product show
\[
 |\det V|=\prod_{j<k}2\left|\sin\frac{39(\omega_k-\omega_j)}2\right|
 >1400.
\]
Hadamard's inequality bounds every $7\times7$ cofactor by
$7^{7/2}<343(8/3)$, so
\[
 \|V^{-1}\|_\infty<\frac{2744}{525},\qquad
 \|\widehat w-w\|_\infty<\frac{43904}{2625}\,e,
\]
where $e$ bounds the total error of each pointer probability.  The latter
includes both quadratures and removal of the known elastic atom.  Finite
precision evaluation of an implemented inverse would need its own error
allowance.

A sufficient prospective budget is systematic probability error
$1/65536$ and sampling error $1/32768$ at each of the sixteen settings.
With $9\cdot2^{29}$ independent trials per setting, the simultaneous
failure probability is less than $32(3/8)^9<1/100$ and the weight error is
less than $343/448000$.  Twice that error is smaller than the least
certified weight, so every line is separated from zero.  The total is
$77\,309\,411\,328$ trials and the largest model time is 273.
This long-time accuracy and sampling schedule is not executed here.
It identifies weights on the supplied support, not an arbitrary unknown
spectrum from finitely many samples.

The finite result is an internally observable protected excitation response
of the admitted simulator, with complete-history and finite-time errors.
It does not require an external particle name or an SI calibration.  It also
does not establish that these neutral modes are single odd fermions, that
the reference is a vacuum ground state, or that finite spectral lines are
relativistic rest masses.  Transported excitations, many-body continuum
dynamics and comparison with nature remain separate constructions.
